% GENERATED FILE. Edit canonical lessons, then run npm run generate:books.
% canonical-source-hash: fnv1a64:7bfbf683eb925d87
% canonical-lessons: ML-C518-bhayam, ML-C518-lajjikkuka, ML-C518-ekantata, ML-C518-kotikkuka, ML-C518-anandam

\chapter{Dread and Longing}
\label{ch:ml-dread-and-longing}
\hlchaptermodality{518}

\begin{chapteropening}
\textbf{By the end of this chapter:} \emph{I can say dread, to feel ashamed, loneliness, to long for and bliss.}

Complete the last lesson of chapter 518: I can say dread, to feel ashamed, loneliness, to long for and bliss.
\end{chapteropening}

\section[\texorpdfstring{\ml{ഭയം} (bhayaṁ)}{ഭയം (bhayaṁ)}]{\ml{ഭയം} (bhayaṁ) — dread; fear}

\label{lesson:ML-C518-bhayam}



\begin{quote}
Before the new one: say the Malayalam for wrath, then the Malayalam for to be upset.
\end{quote}

\subsection*{You'll want to know: \ml{ഭയം}}
\textbf{\ml{ഭയം}} — \emph{bhayaṁ} — \textquotedblleft{}dread; fear\textquotedblright{}.

\begin{quote}
\ml{കുട്ടിക്ക്} \ml{ഇരുട്ടിനെ} \ml{ഭയമായിരുന്നു}.
\end{quote}

\begin{quote}
\emph{kuṭṭikkŭ iruṭṭine bhayamāyirunnu.}
\end{quote}

\begin{quote}
The child was afraid of the dark.
\end{quote}

Written and formal; in speech \ml{പേടി} is more usual. The one afraid takes -\ml{ക്ക്}, the thing feared -\ml{എ}.

\begin{cousinweb}
\textbf{\ml{ഭയം}} (\emph{bhayaṁ}) is borrowed from Sanskrit \emph{bhaya}, \textquotedblleft{}fear, peril\textquotedblright{}, from the root \emph{bhī}, \textquotedblleft{}to fear\textquotedblright{} (CDIAL 9390).
\end{cousinweb}

\begin{grammarlens}[title={What you've built}]
One more word for describing what you lived and felt.
\end{grammarlens}

\subsection*{Your turn}
\begin{itemize}
\raggedright
  \item \emph{Say it:} \emph{bhayaṁ}
  \item \emph{Say it:} \emph{bhayaṁ}, once more
  \item \emph{Say it:} \emph{kuṭṭikkŭ iruṭṭine bhayamāyirunnu.} — The child was afraid of the dark.
  \item \emph{Recall:} say the Malayalam for wrath, then the Malayalam for to be upset, then say \emph{bhayaṁ} again
\end{itemize}

\begin{tcolorbox}[breakable,colback=teal!4,colframe=teal!35!black,title={Before you move on}]
What does \ml{ഭയം} mean? (\textquotedblleft{}Dread; fear\textquotedblright{}.) Say it once more.
\end{tcolorbox}
\section[\texorpdfstring{\ml{ലജ്ജിക്കുക} (lajjikkuka)}{ലജ്ജിക്കുക (lajjikkuka)}]{\ml{ലജ്ജിക്കുക} (lajjikkuka) — to feel ashamed}

\label{lesson:ML-C518-lajjikkuka}



\begin{quote}
Before the new one: say the Malayalam for to be upset, then the Malayalam for dread.
\end{quote}

\subsection*{You'll want to know: \ml{ലജ്ജിക്കുക}}
\textbf{\ml{ലജ്ജിക്കുക}} — \emph{lajjikkuka} — \textquotedblleft{}to feel ashamed\textquotedblright{}.

\begin{quote}
\ml{തെറ്റ്} \ml{ചെയ്തതിൽ} \ml{അവൻ} \ml{ലജ്ജിച്ചു}.
\end{quote}

\begin{quote}
\emph{teṟṟŭ ceytatil avan lajjiccu.}
\end{quote}

\begin{quote}
He felt ashamed of doing wrong.
\end{quote}

Past: \ml{ലജ്ജിച്ചു}, \emph{lajjiccu}.

The noun is \ml{ലജ്ജ}, shame; what one is ashamed of takes -\ml{ഇൽ}.

\begin{cousinweb}
\textbf{\ml{ലജ്ജിക്കുക}} (\emph{lajjikkuka}) is built from Sanskrit \emph{lajjā}, \textquotedblleft{}shame, modesty\textquotedblright{} (CDIAL 10910), with the Malayalam verb ending -\ml{ിക്കുക}.
\end{cousinweb}

\begin{grammarlens}[title={What you've built}]
One more word for describing what you lived and felt.
\end{grammarlens}

\subsection*{Your turn}
\begin{itemize}
\raggedright
  \item \emph{Say it:} \emph{lajjikkuka}
  \item \emph{Say it:} \emph{lajjikkuka}, once more
  \item \emph{Say it:} \emph{teṟṟŭ ceytatil avan lajjiccu.} — He felt ashamed of doing wrong.
  \item \emph{Recall:} say the Malayalam for to be upset, then the Malayalam for dread, then say \emph{lajjikkuka} again
\end{itemize}

\begin{tcolorbox}[breakable,colback=teal!4,colframe=teal!35!black,title={Before you move on}]
What does \ml{ലജ്ജിക്കുക} mean? (\textquotedblleft{}To feel ashamed\textquotedblright{}.) Say it once more.
\end{tcolorbox}
\section[\texorpdfstring{\ml{ഏകാന്തത} (ēkāntata)}{ഏകാന്തത (ēkāntata)}]{\ml{ഏകാന്തത} (ēkāntata) — loneliness; solitude}

\label{lesson:ML-C518-ekantata}



\begin{quote}
Before the new one: say the Malayalam for dread, then the Malayalam for to feel ashamed.
\end{quote}

\subsection*{You'll want to know: \ml{ഏകാന്തത}}
\textbf{\ml{ഏകാന്തത}} — \emph{ēkāntata} — \textquotedblleft{}loneliness; solitude\textquotedblright{}.

\begin{quote}
\ml{നഗരത്തിൽ} \ml{മുത്തശ്ശിക്ക്} \ml{ഏകാന്തത} \ml{തോന്നി}.
\end{quote}

\begin{quote}
\emph{nagarattil muttaśśikkŭ ēkāntata tōnni.}
\end{quote}

\begin{quote}
In the city, grandmother felt lonely.
\end{quote}

Built on \ml{ഏകാന്തം}, a lonely place, with -\ml{ത}, which makes a noun of a quality.

\begin{cousinweb}
\textbf{\ml{ഏകാന്തത}} (\emph{ēkāntata}) is borrowed from Sanskrit \emph{ekānta}, \textquotedblleft{}a lonely or retired place\textquotedblright{} (Monier-Williams), with the noun ending \emph{-tā}. The Sanskrit word joins \emph{eka}, \textquotedblleft{}one\textquotedblright{}, and \emph{anta}, \textquotedblleft{}end, side\textquotedblright{}.
\end{cousinweb}

\begin{grammarlens}[title={What you've built}]
One more word for describing what you lived and felt.
\end{grammarlens}

\subsection*{Your turn}
\begin{itemize}
\raggedright
  \item \emph{Say it:} \emph{ēkāntata}
  \item \emph{Say it:} \emph{ēkāntata}, once more
  \item \emph{Say it:} \emph{nagarattil muttaśśikkŭ ēkāntata tōnni.} — In the city, grandmother felt lonely.
  \item \emph{Recall:} say the Malayalam for dread, then the Malayalam for to feel ashamed, then say \emph{ēkāntata} again
\end{itemize}

\begin{tcolorbox}[breakable,colback=teal!4,colframe=teal!35!black,title={Before you move on}]
What does \ml{ഏകാന്തത} mean? (\textquotedblleft{}Loneliness; solitude\textquotedblright{}.) Say it once more.
\end{tcolorbox}
\section[\texorpdfstring{\ml{കൊതിക്കുക} (kotikkuka)}{കൊതിക്കുക (kotikkuka)}]{\ml{കൊതിക്കുക} (kotikkuka) — to long for; to crave}

\label{lesson:ML-C518-kotikkuka}



\begin{quote}
Before the new one: say the Malayalam for to feel ashamed, then the Malayalam for loneliness.
\end{quote}

\subsection*{You'll want to know: \ml{കൊതിക്കുക}}
\textbf{\ml{കൊതിക്കുക}} — \emph{kotikkuka} — \textquotedblleft{}to long for; to crave\textquotedblright{}.

\begin{quote}
\ml{നാട്ടിലെ} \ml{ഭക്ഷണം} \ml{കഴിക്കാൻ} \ml{ഞാൻ} \ml{കൊതിച്ചു}.
\end{quote}

\begin{quote}
\emph{nāṭṭile bhakṣaṇaṁ kaḻikkān ñān koticcu.}
\end{quote}

\begin{quote}
I longed to eat the food of home.
\end{quote}

Past: \ml{കൊതിച്ചു}, \emph{koticcu}.

What you long for takes -\ml{ആൻ} (a verb) or -\ml{ഇന്} (a noun). The noun is \ml{കൊതി}, a craving.

\begin{cousinweb}
\textbf{\ml{കൊതിക്കുക}} (\emph{kotikkuka}) is native Dravidian: DEDR 2084 lists Malayalam \ml{കൊതി}, \textquotedblleft{}eagerness, greed\textquotedblright{}, and the verb, \textquotedblleft{}to covet, to crave\textquotedblright{} (DEDR 2084).
\end{cousinweb}

\begin{grammarlens}[title={What you've built}]
One more word for describing what you lived and felt.
\end{grammarlens}

\subsection*{Your turn}
\begin{itemize}
\raggedright
  \item \emph{Say it:} \emph{kotikkuka}
  \item \emph{Say it:} \emph{kotikkuka}, once more
  \item \emph{Say it:} \emph{nāṭṭile bhakṣaṇaṁ kaḻikkān ñān koticcu.} — I longed to eat the food of home.
  \item \emph{Recall:} say the Malayalam for to feel ashamed, then the Malayalam for loneliness, then say \emph{kotikkuka} again
\end{itemize}

\begin{tcolorbox}[breakable,colback=teal!4,colframe=teal!35!black,title={Before you move on}]
What does \ml{കൊതിക്കുക} mean? (\textquotedblleft{}To long for; to crave\textquotedblright{}.) Say it once more.
\end{tcolorbox}
\section[\texorpdfstring{\ml{ആനന്ദം} (ānandaṁ)}{ആനന്ദം (ānandaṁ)}]{\ml{ആനന്ദം} (ānandaṁ) — bliss; great joy}

\label{lesson:ML-C518-anandam}



\begin{quote}
Before the new one: say the Malayalam for loneliness, then the Malayalam for to long for.
\end{quote}

\subsection*{You'll want to know: \ml{ആനന്ദം}}
\textbf{\ml{ആനന്ദം}} — \emph{ānandaṁ} — \textquotedblleft{}bliss; great joy\textquotedblright{}.

\begin{quote}
\ml{കുഞ്ഞിനെ} \ml{കണ്ടപ്പോൾ} \ml{മുത്തച്ഛന്റെ} \ml{മുഖത്ത്} \ml{ആനന്ദം} \ml{നിറഞ്ഞു}.
\end{quote}

\begin{quote}
\emph{kuññine kaṇṭappōḷ muttacchanṟe mukhattŭ ānandaṁ niṟaññu.}
\end{quote}

\begin{quote}
When he saw the baby, joy filled grandfather's face.
\end{quote}

Stronger than \ml{സന്തോഷം}; a word of songs and of writing.

\begin{cousinweb}
\textbf{\ml{ആനന്ദം}} (\emph{ānandaṁ}) is borrowed from Sanskrit \emph{ānanda}, \textquotedblleft{}happiness, bliss\textquotedblright{} (CDIAL 1172).
\end{cousinweb}

\begin{grammarlens}[title={What you've built}]
One more word for describing what you lived and felt.
\end{grammarlens}

\subsection*{Your turn}
\begin{itemize}
\raggedright
  \item \emph{Say it:} \emph{ānandaṁ}
  \item \emph{Say it:} \emph{ānandaṁ}, once more
  \item \emph{Say it:} \emph{kuññine kaṇṭappōḷ muttacchanṟe mukhattŭ ānandaṁ niṟaññu.} — When he saw the baby, joy filled grandfather's face.
  \item \emph{Recall:} say the Malayalam for loneliness, then the Malayalam for to long for, then say \emph{ānandaṁ} again
\end{itemize}

\begin{tcolorbox}[breakable,colback=teal!4,colframe=teal!35!black,title={Before you move on}]
What does \ml{ആനന്ദം} mean? (\textquotedblleft{}Bliss; great joy\textquotedblright{}.) Say it once more.
\end{tcolorbox}
